\begin{abstract}
Trained immunity and tolerance change how monocytes and macrophages respond to a second
stimulus, but they are measured by destructive restimulation assays. Whether these memory
states leave a signature that label-free imaging could read has not been tested, because no
public dataset images trained or tolerised cells. We asked how far public data can take the
question. We trained a self-supervised video model of innate-cell behaviour on
\val{pre.t1.clips} clips of migrating human neutrophils and tested it, with strong image and
video baselines, on every public dataset that pairs innate-cell imaging or live-cell dynamics
with a molecular, metabolic, signalling or functional readout, summarised as a modality
coverage map. Imaging features separate microglia activation morphologies, and NF-\(\kappa\)B
reporter dynamics encode stimulus identity, but frozen image features did not predict cell
identity or NAD(P)H lifetime beyond simple intensity statistics, and a general-purpose
video model outperformed the cell-specific one. Memory states themselves are paired with
imaging in no public resource. From the restimulation transcriptomes that do exist, after
sample-identity checks, we find that $\beta$-glucan-trained human macrophages have lower TNF
mRNA than untrained cells 4\,h after restimulation in two independent datasets (six of six
donors; \val{pc.212.tnf.d7.mean} and \val{pc.141.TNF.mean} log2), although the same studies
report higher TNF-$\alpha$ protein at 24\,h. A single readout time can therefore show
training as either enhanced or reduced, which shapes the imaging study we propose to test
innate immune memory directly.
\end{abstract}
